% !TEX root = ../main.tex
\chapter{Questions and answers}\label{ch:faq}
Short answers to questions people ask while using the app, each pointing to where the full
answer lives. The last section is for when something goes wrong.

\section{How do I\ldots?}\label{sec:faq-how}

\faq{How do I look at one channel of a recording?}
Open \apppage{explore.corpus}, pick the recording, and open a channel from the map; or go
straight to \apppage{explore.signal}. The trace is drawn at its real scale, with any
detections and annotations over it.

\faq{How do I run an analysis on part of a recording?}
In \apppage{explore.signal}, drag across the trace to mark a \term{span}{span}, then press
\ui{Send span to Analyse}. In \apppage{analyse.chain}, add blocks with \ui{Insert a stage} or
start from a saved \term{template}{template} with \ui{Import}, then press \ui{Run chain}.
\Cref{ch:first-session} walks through this once from start to finish.

\faq{How do I find spikes and measure the gaps between them?}
Detect them with \blockref{detection.spike_v1} or the \code{dehshibi_spikes} template, then
add \blockref{interrogation.intervals} (and, optionally, \blockref{interrogation.trains}).
\Cref{sec:workflows-spikes} gives the steps and the pitfalls.

\faq{How do I run the same chain on every channel?}
Save it with \ui{Save template}, then in \apppage{discovery.runs} choose the scope (recording,
channels, section) and \ui{Apply template}. Each channel gets its own run and its own
null. Discovery offers detection templates only: a chain that ends in an interrogation block
such as Intervals is not listed.

\faq{How do I check whether a result is more than chance?}
In Analyse, turn on the toolbar's \ui{surrogate null} and run again; the footer then reads
what was detected against what the null expects. Discovery runs always carry a null. See
\cref{app:nulls}.

\faq{How do I search for more examples of a shape I have found?}
Promote it in Review with \ui{S}, then press \ui{Seed search in Discovery}. The search, its
null and how to judge the matches are in \cref{sec:workflows-rq2}.

\faq{How do I judge what a chain found?}
Send its detections to a \term{review-queue}{Review queue} (\ui{Send \(n\) unjudged to Review}
in Discovery, or \ui{Pass \(n\) to Review} on an Analyse block), open \apppage{review.inspector},
and press \ui{I} (interesting), \ui{N} (not interesting), \ui{A} (artifact) or \ui{U}
(unsure). \ui{Ctrl Z} undoes the last verdict.

\faq{How do I tell the app what unit a recording is stored in?}
In \apppage{settings.datasets}, set the recording's unit. Blocks that print millivolts read
it; see ``Why does the same setting find nothing on one recording?'' below.

\faq{How do I send heavy work to the cluster?}
When a run's estimate is over the local limit, it does not start on this machine; it becomes
a cluster job instead. In Analyse, pressing \ui{Run} on such a chain writes the SLURM job for
you; a seed search in Discovery offers \ui{Create SLURM script}. The job's script and inputs
are written to a folder; you copy them to the cluster, submit them, and copy the results back
yourself. \apppage{jobs.all} lists every such job and imports a returned result. The limits
are set in \apppage{settings.compute-hpc}; with nothing saved they are 15 minutes in Analyse,
20 for a Discovery template run and 10 for a seed search. \Cref{sec:jobs-sending} has the
details.

\faq{How do I save a chain to use again?}
Press \ui{Save template} on the Analyse toolbar and name it. Library › Templates lists every
saved template, its versions and its runs. A saved template's recipe records every block and
setting, so a later run of it is the same experiment.

\section{Why does\ldots?}\label{sec:faq-why}

\faq{Why does a verdict never change a detection?}
On purpose. A \term{detection}{detection} is what a machine found; an
\term{annotation}{annotation} or an \term{adjudication}{adjudication} is what a person said.
They live in separate tables, and no code may write a human verdict into a machine row or the
reverse (\code{CLAUDE.md}, rule~5). That way you can always ask ``what did the detector find?''
and ``what did people think of it?'' separately, and compare them. Your verdict is stored
beside the detection, never over it.

\faq{Why does the same setting find nothing on one recording and too much on another?}
Some recordings are stored in volts and some in millivolts. A block that is told the
recording's unit (its block entry says ``Unit-aware'') converts for you. Every other block
works in the stored unit, so a threshold of 0.01 means 0.01\,mV on one recording and 10\,mV on
another. \blockref{detection.spike_v1}, whose threshold is in millivolts, is a unit-aware
one; its default still suits only one recording's noise level, so on a noisier recording it
finds too much. An undeclared unit is read as volts, and blocks that print millivolts say so
in amber.

\faq{Why does the same setting mean a different time on another recording?}
Many settings count samples, not seconds. A window of 60 samples is a minute at 1\,Hz and six
seconds at 10\,Hz. Each setting's unit is in its block's settings table in \cref{app:blocks}.

\faq{Why does Analyse not draw a null every time?}
Tuning a chain is meant to take seconds, and a null runs the whole chain again 20 times. So
in Analyse the \ui{surrogate null} toggle is off until you turn it on. Discovery, where claims
are made, always draws one.

\faq{Why can I not open one of the recordings?}
One recording, \code{M4_aug_concat_fs1.mat}, is \term{held-out}{held out}: kept unseen so that
a model can be tested once, at the end, on data nobody tuned anything on. Every page refuses it
until it is unlocked in \apppage{settings.datasets}, which is meant to happen once, after the
project's freeze. A tool that peeked at it would make that final test worthless.

\faq{Why does the null sometimes find more than the real signal?}
Because the method, at those settings, fires on noise of the recording's own character at
least as often as on the recording. That is a finding, not a fault: the result cannot be told
from chance. Change the settings, or the method, and run the null again. The band runs in
\cref{sec:workflows-rq4} are an example.

\faq{Why does a seed search use 200 draws when a template run uses 20?}
A seed search's draw is one fast search over a surrogate; a template run's draw is the whole
chain run again. Both counts are set in \apppage{settings.nulls}, and every page prints the
number actually drawn.

\faq{Why does a run I already made come back instantly?}
Every run is keyed by its \term{recipe}{recipe}: the recording, the span, and every block with
every setting. Run the same recipe again and the app reuses the stored result instead of
recomputing. Change any setting and it is a new run.

\faq{Why does the Library hide some members?}
By default it shows only members larger than the dataset's noise floor (set per dataset in
Settings › Datasets), so that families and their recurrence are counted over real events.
\emph{show sub-floor} brings the others back; nothing is deleted.

\section{What does\ldots{} mean?}\label{sec:faq-what}

\faq{What does ``× null'' mean?}
The real count divided by the average count over the null's draws. ``3.9× null'' means the
real signal gave almost four times what noise of the same kind gave on average. It is a ratio,
not a p-value. \Cref{sec:nulls-reading} explains how to read it.

\faq{What does ``null expects'' mean?}
The average number of detections the same chain made on the surrogate copies of that channel.
Beside it the page prints how many draws that average is over.

\faq{What does ``not comparable'' mean on a comparison?}
A stretch that no run covered, or that no person reviewed. It is never counted as a
disagreement between people and the machine: silence is not a verdict
(\cref{sec:workflows-rq5}).

\faq{What do the five verdict words mean?}
\emph{interesting}, \emph{not\_interesting}, \emph{artifact} (a fault of the recording, not
the organism), \emph{unsure}, and \emph{seed} (an example worth searching for). The list is
fixed. A verdict that fits none of them is recorded as the nearest one with a note.

\faq{What does ``precision'' mean on the scoreboard, and why are there two?}
The share of a run's detections that people accepted. Human labels come in two shapes,
fixed-length review windows and spans drawn around events, and a detection is matched to each
by a different rule. So the scoreboard prints two figures, each with its rule beside it.

\faq{What does ``too short to tell'' mean?}
In the cross-channel test, a member under 30 samples is too short for its correlation with
another channel to mean anything, so it is neither counted as an artifact nor as a recurrence.

\faq{What does ``experimental'' mean on a block?}
That the block was tried on real recordings and did not beat its null there. It is kept so it
can be applied and tested again, and it says so at the start of its description.

\faq{What does a ``demo'' or ``preview'' label mean?}
That the page or panel is showing example data rather than your recordings. The page section
for each part of the app in this manual says which parts are live.

\section{Something went wrong}\label{sec:faq-wrong}

\faq{A red error card appeared. What does it mean?}
The app is built to fail loudly. When part of a page cannot be drawn, it shows a red card with
the error, rather than an empty panel that looks like ``nothing found''. Read the message: it
usually names the block, the setting or the missing data. If you report it, include the
card's text and what you clicked.

\faq{The page says the server returned an error (500). What now?}
A 500 means the Python side raised an error, and the app passes it to the page with its
traceback on purpose. It is a bug or a bad input, never a quiet empty result. Note what you
did, copy the message, and report it. A refusal with a reason (for example the held-out
recording, which is refused with status 423) is not a bug: the message says why.

\faq{A block says my span is too long.}
Some blocks have a ceiling on the span they will take, because their memory grows faster than
the span (the image encodings and the matrix-profile blocks, for example). The Block page
names the stage and its ceiling in samples. Shorten the span, or send the chain to the
cluster (\ui{Create SLURM script}). The ceiling for each block is in its entry in
\cref{app:blocks}.

\faq{A run is very slow.}
Look at the estimate before you press run: pages that can estimate a run show one, including
the cost of its null draws. A seed search over a whole recording, a matrix profile over days,
or a null with many draws can take hours. Shorten the section, lower the draw count in
\apppage{settings.nulls} while exploring, or send the run to the cluster. A Discovery run whose job dies
with the server is marked \emph{failed} with its reason and can be run again.

\faq{A run says ``failed''. Did anything get saved?}
The run's row is kept with the stage it failed at and the error; \ui{View log} shows the
whole log. \unverified{Stages that finished before the failure are cached, so a re-run after a
fix starts from the failed stage.}

\faq{The count changed when I ran the same chain again.}
A run of an unchanged recipe is reused, so a change means something in the recipe changed: a
setting, the span, or the recording's declared unit. Some older stored results were made
before a block was given a fixed random seed; the research file names those and says to re-run
them once with \emph{force}.

\faq{A recording shows gaps, or a block warns about missing samples.}
Missing samples (NaN) are kept as gaps, not filled. Blocks differ in how they treat them: the
drop detector and the symbolic encoders skip a gap and mark it, while others leave every
window that touches a gap unscored. A block that can see the gap before it runs shows a
\emph{Missing samples} warning in amber; its entry in \cref{app:blocks} says what it does.

\faq{The numbers on a page differ from a number in this manual.}
Trust the page. The manual describes the code at one commit (its title page names it), and
the measurements quoted here come from the research files on a given date. If the page and
its own printed rule disagree, that is worth reporting.
